\documentclass[11pt]{article}
\usepackage[margin=1in]{geometry}
\usepackage[T1]{fontenc}
\usepackage{amsmath,amssymb,amsthm,mathtools}
\usepackage{microtype}
\usepackage{enumitem}
\usepackage{xcolor}
\usepackage[colorlinks=true,linkcolor=blue!55!black,
  citecolor=blue!55!black,urlcolor=blue!55!black]{hyperref}
\numberwithin{equation}{section}
\newtheorem{theorem}{Theorem}[section]
\newtheorem{lemma}[theorem]{Lemma}
\newtheorem{proposition}[theorem]{Proposition}
\newtheorem{corollary}[theorem]{Corollary}
\newtheorem{conjecture}[theorem]{Conjecture}
\theoremstyle{definition}
\newtheorem{example}[theorem]{Example}
\newtheorem{remark}[theorem]{Remark}
\newcommand{\cO}{\mathcal O}
\newcommand{\cZ}{\mathcal Z}

\title{An Erd\H{o}s--S\'os-type theorem for antidirected trees}
\author{Louis DeBiasio}
\date{September 6, 2026}
\hypersetup{pdfauthor={Louis DeBiasio}, pdftitle={An Erd\H{o}s--S\'os theorem for antidirected trees}}

\begin{document}
\maketitle

\begin{abstract}
We prove that every digraph $D$ with more than $(t-2)|V(D)|$ arcs contains
every antidirected tree on $t\geq2$ vertices, establishing the conjecture of
Addario-Berry, Havet, Linhares Sales, Reed, and Thomass\'e.  The proof
adapts the permutation-counting argument of the recent Erd\H{o}s--S\'os
proof obtained in the FrontierMath Erd\H{o}s benchmark.  The result
also gives Sumner's tournament bound for antidirected trees.
An accompanying Lean~4 formalization proves the density theorem.
\end{abstract}

\section*{Acknowledgments}
ChatGPT~6 Astra was used to develop the argument, write the exposition,
and produce the accompanying Lean formalization.  The author's role
was minimal, consisting mostly of asking questions towards trying to understand
the recent Erd\H{o}s--S\'os proof and its possible extensions.

\section{Introduction}
\label{sec:introduction}

An \emph{oriented tree} is obtained by assigning a direction to every
edge of a tree.  It is \emph{antidirected} if every vertex is a source
or a sink; equivalently, no vertex has both positive indegree and
positive outdegree.  We prove an analogue of the Erd\H{o}s--S\'os
theorem for these trees in digraphs.

Throughout the paper, a digraph is finite and loopless, but it may
contain both opposite arcs between a pair of vertices.  There is at most
one arc in each ordered direction.  Copies are injective and need not be
induced, and every target tree has $t\geq2$ vertices.  We write $e(D)$ for
the number of arcs of $D$.

We prove the following density statement, conjectured
by Addario-Berry, Havet, Linhares Sales, Reed, and Thomass\'e in 2013
\cite{AddarioBerryHavetLinharesSalesReedThomasse2013}.

\begin{theorem}
\label{thm:antidirected-density}
Let $T$ be an antidirected tree on $t\geq2$ vertices.  Every digraph $D$ satisfying
\[
        e(D)>(t-2)|V(D)|
\]
contains a copy of $T$.
\end{theorem}

\subsection{Related work}

Several special cases were known.  Addario-Berry et al. proved the
conjecture for antidirected trees of diameter at most three
\cite{AddarioBerryHavetLinharesSalesReedThomasse2013}.  Stein and
Trujillo-Negrete later proved it for every antidirected caterpillar, and for
every antidirected tree when the host excludes three specified orientations
of $K_{2,\lceil (t-1)/12\rceil}$
\cite{SteinTrujilloNegrete2025}.  Stein and Z\'arate-Guer\'en also obtained
an asymptotic version for balanced antidirected trees of bounded maximum
degree and linear order in large oriented graphs
\cite{SteinZarateGueren2024}.

\subsection{Connection with the Erd\H{o}s--S\'os proof}

Our proof is based on the new proof of the Erd\H{o}s--S\'os conjecture
obtained in the FrontierMath Erd\H{o}s benchmark and recorded in
\cite{Erdos548Repository}.\footnote{The repository is maintained by
Tom Adamczewski.  Its provenance account credits the graph proof to
GPT-6 Astra, working autonomously in the FrontierMath Erd\H{o}s
benchmark of Adamczewski and Bloom.}
We adapt its counting of marked vertex permutations and induction on
rooted trees; we do not claim an independent origin for this method.
The directed argument distinguishes arcs leaving the first vertex of
a permutation from arcs entering it.  We check that the two operations
in the graph proof, adding a leaf and joining branches at a root,
respect the required arc directions.  The graph theorem itself is not
used as a black box.

Conversely, our theorem includes the Erd\H{o}s--S\'os assertion for
graphs as a special case.  Given a
graph $G$, form the symmetric digraph $\overleftrightarrow{G}$ by replacing
each edge with its two opposite orientations.  Given a tree $U$, choose its
bipartition and orient every edge from one class to the other, obtaining an
antidirected tree $T$.  Then
\[
 e(\overleftrightarrow{G})=2e(G),
 \qquad
 \overleftrightarrow{G}\text{ contains }T
 \quad\Longleftrightarrow\quad
 G\text{ contains }U.
\]
Thus Theorem~\ref{thm:antidirected-density}, applied to
$\overleftrightarrow{G}$, specializes precisely to the Erd\H{o}s--S\'os
assertion that $2e(G)>(t-2)|V(G)|$ forces a copy of $U$.

\subsection{Outline of the proof}

For each permutation $(v_1,\ldots,v_N)$ of the host vertices, we mark
position $i\geq2$ according to whether $v_1\to v_i$ or $v_i\to v_1$.
We count pairs consisting of a permutation and a marked position for
which the first $i$ vertices contain a copy of the target tree, with
its root mapped to $v_1$.  Each pair is counted once, regardless of
how many such copies it contains.

The induction has two cases.  If the root is a leaf, we start with a
copy of the smaller tree lying before the current marked position.
The vertex at that position is then available for the new leaf.
If the root is not a leaf, we split the tree into two smaller trees
meeting only at the root.  Reordering the permutation gives a counting
inequality in which the two copies use disjoint vertices outside that
root.  The direction of the marking arc is unchanged when branches
are joined, and is reversed relative to the root when a leaf becomes
the new root.  In an antidirected tree, neighbors have opposite roles
as sources and sinks, while branches at a vertex share the same role.
These facts allow the induction to proceed.

Section~\ref{sec:marks} defines the marked pairs and computes their
number.  Section~\ref{sec:path-warmup} gives a separate proof for
antidirected paths, illustrating the sign-switching reversal.
Section~\ref{sec:signed-injections} proves the two counting inequalities
for rooted trees, and Section~\ref{sec:density-proof} completes the
induction.  The tournament consequence and the Lean formalization
are given in Sections~\ref{sec:tournaments} and~\ref{sec:lean}.

\section{Marked permutations}
\label{sec:marks}

Let $D$ be a digraph on $N\geq1$ vertices, and let
\[
        \pi=(v_1,v_2,\ldots,v_N)
\]
be an ordering of its vertices.  A position $i\geq2$ is \emph{positive} if
$v_1\to v_i$ and \emph{negative} if $v_i\to v_1$.  Write
\[
\begin{aligned}
 \cO_D^+
   &=\{(\pi,i):2\leq i\leq N,\ v_1\to v_i\},\\
 \cO_D^-
   &=\{(\pi,i):2\leq i\leq N,\ v_i\to v_1\},
\end{aligned}
\]
and put $M^+=|\cO_D^+|$ and $M^-=|\cO_D^-|$.
A position may have both signs when both opposite arcs are present, and it
has neither sign when the two vertices are nonadjacent.

\begin{proposition}[Basic counts]
\label{prop:digraph-normalization}
For every digraph $D$ on $N\geq1$ vertices,
\begin{equation}
        M^+=M^-=e(D)(N-1)!.
        \label{eq:digraph-normalization}
\end{equation}
\end{proposition}

\begin{proof}
If the first vertex is $x$, every ordering of the remaining vertices has
$d_D^+(x)$ positive positions.  Hence
\[
 M^+=(N-1)!\sum_{x\in V(D)}d_D^+(x)
     =e(D)(N-1)!.
\]
The negative count is identical, using indegrees instead of outdegrees.
\end{proof}

For a pair $(\pi,i)$, its \emph{prefix} is the vertex set
$\{v_1,\ldots,v_i\}$.  Let $(H,h)$ be a rooted oriented graph.
For $\varepsilon\in\{+,-\}$, define $\cO_D^\varepsilon(H,h)$ to be
the set of pairs in $\cO_D^\varepsilon$ whose prefix contains a copy
of $H$ mapping $h$ to $v_1$.  We say that such a pair \emph{supports}
$(H,h)$.  The copy need not follow the order of the permutation,
and each pair is counted only once.  When the host is fixed, we write
$\cO^\varepsilon(H,h)$.  The notation $-\varepsilon$ denotes the sign
opposite to $\varepsilon$.

\section{Warm-up: antidirected paths}
\label{sec:path-warmup}

The density bound in Theorem~\ref{thm:path-density} is already known:
it is a special case of the antidirected-caterpillar theorem of Stein
and Trujillo-Negrete \cite{SteinTrujilloNegrete2025}.  We give a
separate permutation-counting proof to illustrate the sign-switching
mechanism.

For paths rooted at an endpoint, we need only one operation: extend
an earlier path by the vertex at the current marked position and make
that vertex the new root.  The root changes from a source to a sink,
or conversely.  Prefix reversal changes the marking sign in exactly
the same way.  We prove the path case separately to isolate this
mechanism before introducing branch joining.

For $k\geq1$ and $\varepsilon\in\{+,-\}$, let $P_k^\varepsilon$ be
the antidirected path with $k$ arcs, rooted at an endpoint that is a
source if $\varepsilon=+$ and a sink if $\varepsilon=-$.  Thus the
vertex sequence starts $x_0\to x_1\leftarrow x_2\to\cdots$ in the
positive case and $x_0\leftarrow x_1\to x_2\leftarrow\cdots$ in the
negative case.  Let $A_k^\varepsilon$ count the marked states of sign
$\varepsilon$ whose prefixes contain a copy of $P_k^\varepsilon$
rooted at the first vertex.  These are existence counts, not counts
of embeddings.  Every mark of sign $\varepsilon$ supplies the
required one-arc path, so
\begin{equation}
 A_1^\varepsilon=M^\varepsilon=e(D)(N-1)!.
 \label{eq:path-base}
\end{equation}

\begin{lemma}[Extending an antidirected path]
\label{lem:path-extension}
For every $k\geq1$ and $\varepsilon\in\{+,-\}$,
\begin{equation}
 A_k^\varepsilon\leq A_{k+1}^{-\varepsilon}+N!.
 \label{eq:path-extension}
\end{equation}
\end{lemma}

\begin{proof}
For each ordering contributing to $A_k^\varepsilon$, discard its
first supporting marked state $(\pi,i_0)$.  There are at most $N!$
discarded states.  At any remaining state $(\pi,i)$, the earlier
prefix contains $P_k^\varepsilon$ rooted at $v_1$ and avoids $v_i$.
If $\varepsilon=+$, the marking arc is $v_1\to v_i$; adjoining it
extends the path with a new sink endpoint $v_i$.  If
$\varepsilon=-$, the marking arc is $v_i\to v_1$ and the new
endpoint is a source.  In both cases the extended path, rooted at
$v_i$, is $P_{k+1}^{-\varepsilon}$.

Reverse the prefix through position $i$:
\begin{equation}
 (\pi,i)\longmapsto
 \bigl((v_i,v_{i-1},\ldots,v_1,v_{i+1},\ldots,v_N),i\bigr).
 \label{eq:path-reversal}
\end{equation}
The prefix vertex set is unchanged, the first vertex is now $v_i$,
and the marking arc has sign $-\varepsilon$ relative to that vertex.
The image is therefore counted by $A_{k+1}^{-\varepsilon}$.
The index $i$ is retained and a second reversal recovers the source,
so the map is injective.  Adding back the discarded states proves
\eqref{eq:path-extension}.
\end{proof}

The reversal changes the ordering, not the directions of any arcs.
Nor must the supported path follow the permutation order: only its
root and its containing prefix are prescribed.

\begin{example}[A source endpoint becomes a sink]
Let $D$ have vertices $a,b,c,d$ and arcs
\[
 a\to b,\qquad c\to b,\qquad a\to c,\qquad a\to d.
\]
For $\pi=(a,b,c,d)$, position $3$ is a positive mark supporting the
two-arc path $a\to b\leftarrow c$ rooted at $a$.  The marking arc
$a\to c$ is not part of that path.  At the later positive mark $4$,
the earlier copy avoids $d$, so adjoining $a\to d$ gives
\[
 d\leftarrow a\to b\leftarrow c.
\]
Reversal sends $(\pi,4)$ to $((d,c,b,a),4)$.  The new state is
negative, with marking arc $a\to d$, and supports the displayed
three-arc path rooted at its sink endpoint $d$.
\end{example}

\begin{theorem}[Antidirected-path density bound]
\label{thm:path-density}
Let $P$ be any antidirected path on $t\geq2$ vertices.  If a digraph
$D$ contains no copy of $P$, then
\[
 e(D)\leq(t-2)|V(D)|.
\]
\end{theorem}

\begin{proof}
The empty host is immediate, so put $N=|V(D)|\geq1$.  We claim that
for every $k\geq1$ and each sign $\varepsilon$,
\begin{equation}
 e(D)(N-1)!\leq A_k^\varepsilon+(k-1)N!.
 \label{eq:path-counting}
\end{equation}
The case $k=1$ is \eqref{eq:path-base}.  For $k>1$, apply the
inductive inequality with sign $-\varepsilon$, followed by
Lemma~\ref{lem:path-extension}, to obtain
\[
 e(D)(N-1)!
 \leq A_{k-1}^{-\varepsilon}+(k-2)N!
 \leq A_k^\varepsilon+(k-1)N!.
\]
This proves the claim.  Choose an endpoint of $P$ and let
$\varepsilon$ be its source or sink sign.  Since $D$ is $P$-free,
$A_{t-1}^\varepsilon=0$.  Taking $k=t-1$ in
\eqref{eq:path-counting} and dividing by $(N-1)!$ proves the bound.
\end{proof}

The coefficient is sharp: a complete symmetric digraph on $t-1$
vertices has $(t-1)(t-2)$ arcs and cannot contain $P$.
The same sign-switching reversal will add a leaf to an arbitrary
rooted oriented tree in Lemma~\ref{lem:signed-leaf}.  The extra
operation needed for general antidirected trees is joining branches
at a common root, which we develop next.

\section{Joining copies and adding a leaf}
\label{sec:signed-injections}

We first compare the counts for two trees that meet at a single root.
This step does not require the trees to be antidirected.  That
assumption enters when we use both counting inequalities in the induction.

\begin{lemma}[Joining at a root]
\label{lem:signed-branch}
Let $T_1$ and $T_2$ be rooted oriented subtrees of an oriented tree $T$,
all rooted at $r$, such that
\[
 T_1\cup T_2=T,
 \qquad
 V(T_1)\cap V(T_2)=\{r\}.
\]
Then, for either sign $\varepsilon$,
\begin{equation}
 |\cO^\varepsilon(T_1,r)|+|\cO^\varepsilon(T_2,r)|
 \leq M^\varepsilon+N!+|\cO^\varepsilon(T,r)|.
 \label{eq:signed-branch}
\end{equation}
\end{lemma}

\begin{proof}
We also allow the pairs
\[
 \cZ=\{(\pi,1):\pi\text{ is an ordering of }V(D)\},
 \qquad |\cZ|=N!.
\]
We construct an injection
\[
 \cO^\varepsilon(T_1,r)\setminus\cO^\varepsilon(T,r)
 \longrightarrow
 \bigl(\cO_D^\varepsilon\setminus\cO^\varepsilon(T_2,r)\bigr)
 \cup\cZ.
\]
Take a pair $(\pi,i)$ in the domain and let $a\leq i$ be its first
position of sign $\varepsilon$ supporting $(T_1,r)$.  Write
\[
 \pi=(x,A,B,C),\qquad
 A=(v_2,\ldots,v_a),\quad B=(v_{a+1},\ldots,v_i),\quad x=v_1.
\]
Here $A,B,C$ are consecutive lists of vertices; the original prefix
ends after $B$.  Set $b=1+|B|$ and send the pair to
\begin{equation}
 ((x,A,B,C),i)\longmapsto ((x,B,A,C),b).
 \label{eq:signed-rotation}
\end{equation}
If $B$ is empty, the image belongs to $\cZ$.  Otherwise its first
vertex is still $x$ and the last vertex of its prefix is still $v_i$,
so its marking arc has the same sign $\varepsilon$.

The vertices in $x,A$ support a copy of $T_1$ rooted at $x$.
If the image prefix $x,B$ supported $T_2$ rooted at $x$, the two
copies would be disjoint outside $x$ and could be joined to embed
$T$ in the original prefix.  This contradicts the source condition.
Thus the image belongs to the stated set.

For the inverse, read $B$ from the image prefix after $x$.  Starting
immediately after that prefix, recover $A$ as the shortest nonempty
initial list whose last vertex supplies a mark of sign $\varepsilon$
at $x$ and which, together with $x$, supports $T_1$ rooted at $x$.
The original $A$ satisfies these conditions, and no proper initial
segment does, by the minimality of $a$.  This test does not include
$B$.  The remaining entries are $C$, and the original index is
$i=1+|A|+|B|$.  This also recovers the source when $b=1$.
The map is therefore injective.  Since support of $T$ implies support
of $T_1$, taking cardinalities yields
\[
 |\cO^\varepsilon(T_1,r)|-|\cO^\varepsilon(T,r)|
 \leq M^\varepsilon-|\cO^\varepsilon(T_2,r)|+N!,
\]
which is \eqref{eq:signed-branch}.
\end{proof}

We next add a leaf and make it the new root, using the same
construction as in Lemma~\ref{lem:path-extension}, now for an arbitrary
rooted oriented tree.  Reversing the prefix interchanges the first
and marked vertices.  The marking arc itself
does not change, but its sign does: an arc entering the old root
leaves the new root, and conversely.

\begin{lemma}[Adding a leaf]
\label{lem:signed-leaf}
Let $(S,p)$ be a rooted oriented tree.  Add a new leaf $\ell$ adjacent to
$p$, producing $T$, and root $T$ at $\ell$.  Define
\[
 \sigma=
 \begin{cases}
  +,&\ell\to p,\\
  -,&p\to\ell.
 \end{cases}
\]
Then
\begin{equation}
        |\cO^{-\sigma}(S,p)|
        \leq |\cO^\sigma(T,\ell)|+N!.
        \label{eq:signed-leaf}
\end{equation}
\end{lemma}

\begin{proof}
Let $\cZ$ again be the $N!$ pairs $(\pi,1)$.  We construct an injection
\[
 \Psi:\cO^{-\sigma}(S,p)
       \longrightarrow \cO^\sigma(T,\ell)\cup\cZ.
\]
For each ordering $\pi$, send its first pair in
$\cO^{-\sigma}(S,p)$, if one exists, to $(\pi,1)\in\cZ$.

For any later pair $(\pi,i)$, an earlier prefix contains $S$ rooted at
$v_1$, and therefore does not use $v_i$.  Reverse the first $i$ entries:
\begin{equation}
 (\pi,i)\longmapsto
 ((v_i,v_{i-1},\ldots,v_1,v_{i+1},\ldots,v_N),i).
 \label{eq:signed-reversal}
\end{equation}
If $\sigma=+$, then the source is
negative, so $v_i\to v_1$; this arc represents $\ell\to p$ and is positive
after the reversal.  If $\sigma=-$, then the source is positive, so
$v_1\to v_i$; this arc represents $p\to\ell$ and is negative after the
reversal.  In either case the image supports $(T,\ell)$ and has sign
$\sigma$.  Prefix reversal is an involution for each fixed $i$, and the
images in $\cZ$ are distinct from one another and from the images with
$i\geq2$, so the map is injective.
\end{proof}

\section{Proof of the directed density theorem}
\label{sec:density-proof}

For a vertex $r$ of an antidirected tree $T$, put
\[
 \sigma_T(r)=
 \begin{cases}
  +,&r\text{ is a source},\\
  -,&r\text{ is a sink}.
 \end{cases}
\]
Adjacent vertices have opposite signs.  If we split a tree into branches
at a root, the root remains a source in every branch or a sink in every
branch.  Thus Lemma~\ref{lem:signed-leaf} uses the appropriate sign
at the new root, and Lemma~\ref{lem:signed-branch} uses the same sign
in both smaller trees.

\begin{theorem}[The main counting inequality]
\label{thm:antidirected-master}
Let $T$ be an antidirected tree on $t\geq2$ vertices, rooted at $r$, and
let $D$ be a digraph on $N\geq1$ vertices.  With
$\sigma=\sigma_T(r)$,
\begin{equation}
        M^\sigma
        \leq |\cO^\sigma(T,r)|+(t-2)N!.
        \label{eq:antidirected-master}
\end{equation}
\end{theorem}

\begin{proof}
We induct on $t$.  If $t=2$, the unique arc incident with $r$ points away
from $r$ when $\sigma=+$ and towards $r$ when $\sigma=-$.  Every pair of
sign $\sigma$ therefore supplies the required rooted edge, so
$M^\sigma=|\cO^\sigma(T,r)|$.

Suppose first that $r$ is a leaf, let $p$ be its neighbour, and put
$S=T-r$.  The vertex $p$ has sign $-\sigma$ in $S$.  By induction and
Proposition~\ref{prop:digraph-normalization},
\[
 M^\sigma=M^{-\sigma}
 \leq |\cO^{-\sigma}(S,p)|+(t-3)N!.
\]
Lemma~\ref{lem:signed-leaf} now gives
\[
 M^\sigma
 \leq |\cO^\sigma(T,r)|+(t-2)N!.
\]

Suppose instead that $r$ has degree at least two.  Partition the components
of $T-r$ into two nonempty groups, and let $T_1,T_2$ be the corresponding
rooted subtrees after restoring $r$.  Both have root sign $\sigma$.  If
$t_j=|V(T_j)|$, then $t_1+t_2=t+1$.  Adding the two induction inequalities
and applying Lemma~\ref{lem:signed-branch} gives
\[
\begin{aligned}
 2M^\sigma
   &\leq |\cO^\sigma(T_1,r)|+|\cO^\sigma(T_2,r)|
          +(t_1+t_2-4)N!\\
   &\leq M^\sigma+|\cO^\sigma(T,r)|
          +(t-2)N!.
\end{aligned}
\]
Subtracting $M^\sigma$ proves the result.
\end{proof}

\begin{proof}[Proof of Theorem~\ref{thm:antidirected-density}]
The strict density hypothesis excludes the empty host.  Write
$N=|V(D)|\geq1$, and root $T$ at any vertex $r$.  If $D$ were
$T$-free, then $\cO^\sigma(T,r)$ would be empty.  Theorem
\ref{thm:antidirected-master} and Proposition
\ref{prop:digraph-normalization} would give
\[
        e(D)(N-1)!\leq(t-2)N!,
\]
or equivalently $e(D)\leq(t-2)N$, contrary to the hypothesis.
\end{proof}

\section{A tournament consequence}
\label{sec:tournaments}

Sumner's conjecture asserts that every tournament on $2t-2$ vertices
contains every oriented tree on $t$ vertices; it is known for all
sufficiently large $t$ \cite{KuhnMycroftOsthus2011,SteinSurvey2024}.
Theorem~\ref{thm:antidirected-density} gives the antidirected case
for every $t\geq2$.

\begin{corollary}[Sumner's bound for antidirected trees]
Every tournament on at least $2t-2$ vertices contains every antidirected
tree on $t\geq2$ vertices.
\end{corollary}

\begin{proof}
Restrict to any $N=2t-2$ vertices.  The resulting tournament $D$ has
\[
 e(D)=\frac{N(N-1)}2=N\left(t-\frac32\right)>(t-2)N.
\]
Apply Theorem~\ref{thm:antidirected-density}.
\end{proof}

\section{Lean formalization}
\label{sec:lean}

The accompanying project in \texttt{lean/} proves the counting
inequalities and the density theorem in Lean~4.19.0, using the bundled
\texttt{Std} library.  We describe its definitions and their relation
to the proof.  Instructions for checking the files are given in
\texttt{lean/README.md}.

\subsection{Definitions and theorem statements}

Finite digraphs are represented by lists of vertices and arcs, with no
repeated entries.  Opposite arcs are allowed.  An embedding preserves
arcs and is injective on all target vertices, not just separately on
the sources and the sinks.

The final density theorem assumes that the target has at least two
vertices, is antidirected, and has a connected underlying graph with
no simple cycle.  Connectedness is defined using walks, and a simple
cycle has distinct vertices apart from its repeated endpoint.
These definitions do not refer to the counting argument.
The corresponding declarations are
\begin{itemize}[leftmargin=*]
\item \texttt{Digraph.addarioBerry\_bound\_of\_no\_simple\_cycle},
      giving $e(D)\leq(t-2)|V(D)|$ for a $T$-free host;
\item \texttt{Digraph.addarioBerry\_contains\_of\_no\_simple\_cycle},
      giving Theorem~\ref{thm:antidirected-density}.
\end{itemize}
The older declarations \texttt{Digraph.addarioBerry\_bound} and
\texttt{Digraph.addarioBerry\_contains} use a recursive definition
of the target tree.  Both versions remain available.

\subsection{Relation to the proof}

The counting proof is organized recursively.  It starts with one arc,
adds a new leaf as the root, or joins two trees at their common root.
The code proves that the induction parameter is the number of target
vertices minus two.  It also proves that every antidirected tree
with a connected, acyclic underlying graph can be built by these
operations, with any chosen vertex as its root.  This step is proved
in \texttt{Digraph/TreeBridge.lean}; the final theorems do not require
the user to supply a recursive construction.

The formal proof of Lemma~\ref{lem:signed-leaf} swaps the first and
marked vertices rather than reversing the whole prefix.  Both maps
preserve the prefix vertex set, change the sign of the mark, and are
their own inverses for each fixed marked position.  The proof of
Lemma~\ref{lem:signed-branch} uses the same shortest marked prefix
as the paper.  The path density bound is a special case of the
formalized general theorem; the separate warm-up proof has not been
formalized line by line.  The tournament corollary has not been formalized.

\subsection{Verification}

The project specifies the Lean version and includes a script that
compiles the files and checks which axioms the theorems depend on.
The theorem statements are collected in \texttt{lean/Statements.lean}.
Only \texttt{propext}, \texttt{Classical.choice}, and
\texttt{Quot.sound} are permitted in these dependency checks.

\begin{thebibliography}{99}

\bibitem{AddarioBerryHavetLinharesSalesReedThomasse2013}
L.~Addario-Berry, F.~Havet, C.~Linhares Sales, B.~Reed, and S.~Thomass\'e,
\emph{Oriented trees in digraphs},
Discrete Math. \textbf{313} (2013), 967--974.
\href{https://doi.org/10.1016/j.disc.2013.01.011}
{doi:10.1016/j.disc.2013.01.011}.

\bibitem{KuhnMycroftOsthus2011}
D.~K\"uhn, R.~Mycroft, and D.~Osthus,
\emph{A proof of Sumner's universal tournament conjecture for large
tournaments},
Proc. Lond. Math. Soc. (3) \textbf{102} (2011), 731--766.

\bibitem{SteinTrujilloNegrete2025}
M.~Stein and A.~Trujillo-Negrete,
\emph{Antidirected trees in dense digraphs},
SIAM J. Discrete Math. \textbf{39} (2025), 698--727.
\href{https://doi.org/10.1137/24M165497X}{doi:10.1137/24M165497X}.

\bibitem{SteinSurvey2024}
M.~Stein,
\emph{Oriented trees and paths in digraphs},
in \emph{Surveys in Combinatorics 2024},
London Math. Soc. Lecture Note Ser. \textbf{493},
Cambridge Univ. Press, 2024, 271--295.

\bibitem{SteinZarateGueren2024}
M.~Stein and C.~Z\'arate-Guer\'en,
\emph{Antidirected subgraphs of oriented graphs},
Combin. Probab. Comput. \textbf{33} (2024), 446--466.

\bibitem{Erdos548Repository}
T.~Adamczewski (maintainer),
\emph{Erd\H{o}s problem \#548: Lean proof and provenance}, 2026.
\url{https://github.com/tadamcz/erdos548}.

\end{thebibliography}
\end{document}
